\subsection*{The Relational Character of Kinesis}

A crucial feature of $A_0$ that must be established before we can trace the chain forward is the relational character of \textit{kinesis} itself. Motion is not a property inhering in a single substance; it is a relation between mover and moved, between the active and the passive powers that together constitute the process of change. Aristotle argues at \textit{Physics} III.2, 202a13--20 that ``motion is in the movable. It is the fulfilment of this potentiality by the action of that which has the power of causing motion; and the actuality of that which has the power of causing motion is not other than the actuality of the movable; for it must be the fulfilment of \textit{both}. A thing is capable of causing motion because it \textit{can} do this, it is a mover because it actually \textit{does} it. But it is on the movable that it is capable of acting. Hence there is a single actuality of both alike'' (Aristotle, \textit{Physics}, p.~116). Motion is irreducibly relational, a joint actualization of two capacities---the capacity to act and the capacity to be acted upon---and the \textit{energeia} shared between them is not the consequence of the motion but the motion itself.

Aristotle specifies this dyadic structure in categorial terms in the immediately following section. Motion involves both a \textit{poiētikon} (what can act) and a \textit{pathētikon} (what can be acted upon); the actuality of the agent is \textit{poiesis} and the actuality of the patient is \textit{pathesis}. ``It is necessary that there should be an actuality of the agent and of the patient. The one is agency and the other patiency; and the outcome and end of the one is an action, that of the other a passion'' (\textit{Physics} III.3, 202a21--24). The motion is a single event, yet it has two formal aspects under which it can be described---agency and patiency, \textit{poiesis} and \textit{pathesis}---and these aspects are mutually implicated without being identical. Aristotle clarifies the point with the Thebes--Athens analogy: ``Thebes and Athens are different cities, yet the road from Thebes to Athens and the road from Athens to Thebes are the same'' (\textit{Physics} III.3, 202b13--14). One road, two descriptions. One motion, two formal aspects. The distinction is in account, not in substrate.

The relational character of motion has direct consequences for perception. Perception, as Aristotle defines it in \textit{De Anima}, is a kind of alteration: the sense-organ is moved by the sensible object, and the perceptual act is the joint actualization of the sense-organ's capacity to receive the form and the object's capacity to act upon the organ (Aristotle, \textit{On The Soul (De Anima)}, p.~18). The soul, accordingly, ``must be a substance in the sense of the form of a natural body''---it is the formal principle that organizes the body's material capacities into a unified perceiving whole (Aristotle, \textit{On The Soul (De Anima)}, p.~18). The perceptual event at $A_2$ is not two events conjoined but one event described under two \textit{logoi}, as Aristotle will state explicitly at \textit{De Anima} III.2, 425b26--27. Hence the relational structure of \textit{kinesis} at the $A_0$ level anticipates the relational structure of perception at the $A_1$ level: in both cases, actuality is a joint achievement of agent and patient, mover and moved.

This dyadic architecture, supplied by $A_0$, propagates through every subsequent stage of the chain. $A_0$ establishes that motion is always \textit{pros ti}---always a being-in-relation---and $A_1$ inherits this structure: the sensible object's active potency (\textit{poiesis}) meets the sense-organ's receptive potency (\textit{pathesis}) in a single perceptual actualization. The same schema iterates at the cognitive-motor stages: what the realizable good moves at $M_3 \to M_4$ is \textit{orexis}---the moved mover that is at once agent (of the animal's locomotion) and patient (of the apparent good); what \textit{orexis} in turn moves is the animal, by way of its bodily instrument. The dyadic structure is inherited from $A_0$ to $A_1$ to $A_2$, and it repeats wherever a moved mover appears in the chain.

Burke's grammar of motives is instructive here. If we apply Burke's analysis to \textit{kinesis}, we see that the relational structure of motion is simultaneously a structure of determination: the motion of A by B determines both A and B, assigning each a role (patient and agent) and a trajectory (from potentiality to actuality). Similarly, the perceptual encounter between sense-organ and sensible object determines both: the eye is determined as seeing, the color as seen. A synonym for ``determinations,'' Burke continues, is ``conditions,'' since ``conditions are likewise contextual, as with the conditions of an organism's existence'' (Burke, \textit{A Grammar of Motives}, pp.~214--215). This contextual reading of determination is profoundly Aristotelian: the organism exists within a totality of involvements---a world of sensible objects, media, and temporal horizons---and the conditions of its perceptual existence are given by the relational structures of motion and time that constitute $A_0$. Von Uexküll's concept of the \textit{Umwelt} provides a suggestive parallel, albeit one that must be handled with care. Each organism inhabits its own perceptual world, structured by the specific sense-capacities it possesses and the environmental features to which those capacities are attuned; indeed, as von Uexküll observes, families and professions each have their own \textit{Umwelten}, and ``these two Umwelten must be aligned in such a way that they complement rather than disturb each other'' (von Uexküll, \textit{A Foray Into the Worlds of Animals and Humans with A Theory of Meaning}, pp.~232--234). While von Uexküll's framework is not identical to Aristotle's, the structural affinity is noteworthy: the organism's perceptual world is constituted by the relational structures that link its capacities to its environment, just as $A_0$ constitutes the relational ground upon which $A_1$'s perceptual actualities depend.

A further feature of \textit{kinesis} must be registered for the chain's subsequent stages: motion is the kind of process that can persist even as its originating cause withdraws. The full analysis of this residual structure belongs to the subsequent section on \textit{aisthēsis} ($A_1 \to A_2$), where the residual motion in the sense-organs will be developed as the material condition for the \textit{phantasma} at $A_3$. Here it suffices to note that $A_0$ supplies the ontological possibility of residual motion as a structural feature of \textit{kinesis} itself. Motion, once initiated, does not require the continuous presence of its cause in order to continue; it has a kind of temporal autonomy, grounded in its \textit{ateles} character, that allows residual traces to persist in the moved thing even after the mover has withdrawn.
